%——————————— 中文摘要 ———————————%
\cleardoublepage                           % 保证中文摘要落在右页
\pagestyle{fancy}
\pagenumbering{Roman}
\setcounter{page}{1}
\fancyhead[CE]{\songti\zihao{-5} xxx} % 左页页眉
\fancyhead[CO]{\songti\zihao{-5} 摘要}                             % 右页页眉

\addcontentsline{toc}{section}{摘要}
\section*{摘\quad 要}
%摘要：二字间空两格，黑体三号居中，段前，段后各空一行。

xxx

~\\
%关键字：宋体12磅，行距20磅，段前段后0磅，关键词之间用分号隔开,关键词三个字加粗。
\hspace*{2em}{\heiti \zihao{-4}关键词}：xxx；

% 插入空白页，以保证英文摘要从奇数页开始
\cleardoublepage

%——————————— 英文摘要 ———————————%
\fancyhead[CE]{\songti\zihao{-5} xxx} % 左页页眉
\fancyhead[CO]{\zihao{-5} ABSTRACT}                                % 右页页眉

\addcontentsline{toc}{section}{ABSTRACT}
\titleformat{\section}[block]{\centering\bfseries\thesisenglishfont\fontsize{16pt}{20pt}\selectfont}{\thesection}{1em}{}[]
\section*{ABSTRACT}
%ABSTRCT：Times New Roman 加粗三号，段前段后空一行

xxx
%英文摘要内容：Times New Roman 12磅（即小四号），行距20磅段前段后0磅

~\\
\hspace*{2em}\textbf{Key words}: xxx;\\
%Keywords：Times New Roman 12磅，行距20磅， “key words” 两词加粗
